
\documentclass{article}
\usepackage{colortbl}
\usepackage{makecell}
\usepackage{multirow}
\usepackage{supertabular}

\begin{document}

\newcounter{utterance}

\centering \large Interaction Transcript for game `cladder', experiment `full\_v1.5\_default', episode 127 with Qwen3.6{-}27B.
\vspace{24pt}

{ \footnotesize  \setcounter{utterance}{1}
\setlength{\tabcolsep}{0pt}
\begin{supertabular}{c@{$\;$}|p{.15\linewidth}@{}p{.15\linewidth}p{.15\linewidth}p{.15\linewidth}p{.15\linewidth}p{.15\linewidth}}
    \# & \multicolumn{2}{c}{Player} && \multicolumn{2}{c}{Game Master} \\
    \hline

    \theutterance \stepcounter{utterance}  
    & & & \multicolumn{4}{p{0.6\linewidth}}{
        \cellcolor[rgb]{0.9,0.9,0.9}{
            \makecell[{{p{\linewidth}}}]{
                \texttt{\tiny{[P1$\langle$GM]}}
                \texttt{You are an expert in causal inference. The following question is not a typical commonsense query, but rather a meticulously designed question created by a professor specializing in causal inference, intended to assess the students' mastery of the course content.} \\
\texttt{Imagine a self{-}contained, hypothetical world with only the following conditions, and without any unmentioned factors or causal relationships: CEO has a direct effect on director and manager. Manager has a direct effect on employee. Director has a direct effect on employee. Method 1: We look at how manager correlates with employee case by case according to director. Method 2: We look directly at how manager correlates with employee in general. To understand how manager affects employee, is it more correct to use the Method 1 than Method 2?} \\
\texttt{Start your answer with 'Yes' or 'No', followed by additional reasoning or evidence to support your explanation.} \\
            }
        }
    }
     \\ \\

    \theutterance \stepcounter{utterance}  
    & \multicolumn{4}{p{0.6\linewidth}}{
        \cellcolor[rgb]{0.561,0.737,0.561}{
            \makecell[{{p{\linewidth}}}]{
                \texttt{\tiny{[P1$\rangle$GM]}}
                \texttt{No} \\
\\ 
\texttt{**Reasoning:**} \\
\\ 
\texttt{To determine the correct method for estimating the causal effect of **Manager** on **Employee**, we must analyze the causal structure (Directed Acyclic Graph, or DAG) provided:} \\
\\ 
\texttt{1.  **Causal Structure:**} \\
\texttt{    *   \$CEO \textbackslash{}rightarrow Director\$} \\
\texttt{    *   \$CEO \textbackslash{}rightarrow Manager\$} \\
\texttt{    *   \$Manager \textbackslash{}rightarrow Employee\$} \\
\texttt{    *   \$Director \textbackslash{}rightarrow Employee\$} \\
\\ 
\texttt{2.  **Identify the Path of Interest:**} \\
\texttt{    We want to estimate the direct causal effect of \$Manager\$ on \$Employee\$ (\$Manager \textbackslash{}rightarrow Employee\$).} \\
\\ 
\texttt{3.  **Identify Confounding and Biasing Paths:**} \\
\texttt{    *   There is a backdoor path from \$Manager\$ to \$Employee\$ through \$CEO\$ and \$Director\$:} \\
\texttt{        \$\$Manager \textbackslash{}leftarrow CEO \textbackslash{}rightarrow Director \textbackslash{}rightarrow Employee\$\$} \\
\texttt{    *   This path creates a non{-}causal association (confounding bias) between \$Manager\$ and \$Employee\$ because they share a common cause (\$CEO\$) which also affects \$Employee\$ via \$Director\$.} \\
\texttt{    *   To obtain an unbiased estimate of the causal effect, we must block this backdoor path.} \\
\\ 
\texttt{4.  **Evaluate Method 2 (General Correlation):**} \\
\texttt{    *   Method 2 looks at the marginal correlation between \$Manager\$ and \$Employee\$ without conditioning on any other variables.} \\
\texttt{    *   This leaves the backdoor} \\
            }
        }
    }
    & & \\ \\

    \theutterance \stepcounter{utterance}  
    & & & \multicolumn{2}{p{0.3\linewidth}}{
        \cellcolor[rgb]{0.9,0.9,0.9}{
            \makecell[{{p{\linewidth}}}]{
                \texttt{\tiny{[GM$|$GM]}}
                \texttt{No} \\
            }
        }
    }
    & & \\ \\

    \theutterance \stepcounter{utterance}  
    & & & \multicolumn{2}{p{0.3\linewidth}}{
        \cellcolor[rgb]{0.9,0.9,0.9}{
            \makecell[{{p{\linewidth}}}]{
                \texttt{\tiny{[GM$|$GM]}}
                \texttt{yes} \\
            }
        }
    }
    & & \\ \\

    \theutterance \stepcounter{utterance}  
    & & & \multicolumn{2}{p{0.3\linewidth}}{
        \cellcolor[rgb]{0.9,0.9,0.9}{
            \makecell[{{p{\linewidth}}}]{
                \texttt{\tiny{[GM$|$GM]}}
                \texttt{game\_result = LOSE} \\
            }
        }
    }
    & & \\ \\

\end{supertabular}
}

\end{document}
